\documentclass[11pt]{article}
\usepackage[margin=1.1in]{geometry}
\usepackage{amsmath,amssymb,amsthm}
\usepackage[hidelinks]{hyperref}

\newtheorem{theorem}{Theorem}
\newtheorem{lemma}{Lemma}
\newtheorem{remark}{Remark}
\newcommand{\Lcal}{\mathcal L}
\newcommand{\rhoaus}{\overline{\rho}}

\title{A Doubling-Gauge Variant of Auge's Orbit Theorem}
\author{fa\_banach\_001 / agent\_super\_01}
\date{2026-07-01}

\begin{document}
\maketitle

\section*{Source Question}

Auge's Theorem 4.3 in \emph{Orbits of linear operators and Banach space
geometry}, arXiv:1204.2046, assumes that the modulus of asymptotic uniform
smoothness of \(X\) satisfies
\[
  \rhoaus_X(2t)=O(\rhoaus_X(t)) \qquad (t\downarrow0).
\]
The paper then asks whether this auxiliary assumption is really necessary.
The following observation removes that assumption in the power-gauge regime
used for the source's exponent applications.

\begin{theorem}
\label{thm:gauge}
Let \(X\) be a Banach space, let
\[
  f(t)=\rhoaus_X(t) \qquad (t\ge 0),
\]
and let \(\rho:[0,\infty)\to[0,\infty)\) be nondecreasing, positive on
\((0,\infty)\), and dyadically doubling at zero:
\[
  \rho(2t)=O(\rho(t)) \qquad (t\downarrow0).
\]
Assume
\[
  \lim_{t\downarrow0}\frac{f(t)}{\rho(t)}=0.
\]
Then for every non-nilpotent \(T\in\Lcal(X)\) there exists \(x\in X\) such
that
\[
  \sum_{n=1}^{\infty}
  \rho\!\left(\frac{\|T^n x\|}{\|T^n\|}\right)=\infty .
\]
\end{theorem}

\begin{remark}
Thus Auge's extra doubling hypothesis on \(f=\rhoaus_X\) is not needed when
the comparison gauge is, for example, \(\rho(t)=t^q\). This does not settle
the arbitrary-gauge form of the remark.
\end{remark}

\section*{The Coefficient Lemma}

The proof of Auge's theorem uses the AUS-doubling assumption only to choose
coefficients \(\alpha_i\downarrow0\) such that every dyadic rescaling has a
summable \(f\)-cost while the \(\rho\)-cost still diverges.  We replace that
step by a diagonal lemma.

\begin{lemma}
\label{lem:diag}
Let \(g:(0,\infty)\to(0,\infty)\), and let \(F_k:(0,\infty)\to[0,\infty)\)
\((k\ge0)\) satisfy \(F_k(t)\to0\) and
\[
  \frac{F_k(t)}{g(t)}\longrightarrow0 \qquad (t\downarrow0)
\]
for every fixed \(k\).  Then there is a sequence
\(\alpha_i\downarrow0\) such that
\[
  \sum_i F_k(\alpha_i)<\infty \quad (k=0,1,2,\ldots),
  \qquad
  \sum_i g(\alpha_i)=\infty .
\]
\end{lemma}

\begin{proof}
If \(g\) does not tend to \(0\) at the origin, choose \(\varepsilon>0\) and
\(t_r\downarrow0\) with \(g(t_r)\ge\varepsilon\).  Passing recursively to a
subsequence, we may also arrange \(F_k(t_r)\le2^{-r}\) for every \(k\le r\).
Taking \(\alpha_i=t_i\) gives the result.

Suppose now that \(g(t)\to0\).  For each \(r\ge1\), choose \(t_r\in(0,1/r)\)
so small that \(g(t_r)\le1/2\) and
\[
  F_k(t_r)\le 2^{-r}g(t_r) \qquad (0\le k\le r).
\]
Let \(N_r=\lfloor 1/g(t_r)\rfloor\), and repeat \(t_r\) exactly \(N_r\)
times in the sequence \((\alpha_i)\).  Then the contribution of the \(r\)-th
block to \(\sum_i g(\alpha_i)\) is at least \(1-g(t_r)\ge1/2\), so the
\(g\)-sum diverges.  For fixed \(k\), the tail \(r\ge k\) contributes at most
\[
  \sum_{r\ge k} N_r F_k(t_r)
  \le \sum_{r\ge k} 2^{-r} N_r g(t_r)
  \le \sum_{r\ge k} 2^{-r}<\infty .
\]
The finitely many earlier blocks are harmless.
\end{proof}

\section*{Proof of Theorem \ref{thm:gauge}}

Set \(g(t)=\rho(t/2)\) and \(F_k(t)=f(2^k t)\).  Since \(\rho\) is
dyadically doubling, for every fixed \(k\) there is \(C_k\) such that, for
small \(s\),
\[
  \rho(2^{k+1}s)\le C_k\rho(s).
\]
With \(s=t/2\), the hypothesis \(f=o(\rho)\) gives
\[
  \frac{F_k(t)}{g(t)}
  =
  \frac{f(2^{k+1}s)}{\rho(s)}
  \le
  C_k\,\frac{f(2^{k+1}s)}{\rho(2^{k+1}s)}
  \longrightarrow0.
\]
Also \(F_k(t)\to0\).  Lemma \ref{lem:diag} therefore gives
\(\alpha_i\downarrow0\) such that
\[
  \sum_i f(2^k\alpha_i)<\infty \quad (k=0,1,2,\ldots),
  \qquad
  \sum_i \rho(\alpha_i/2)=\infty .
\]

Now repeat Auge's proof of Theorem 4.3 with this sequence.  The alternative
case in which some finite-codimensional restriction of \(T^n\) is small is
handled exactly as in the source by Proposition 4.1.  In the remaining case,
for every finite-codimensional \(M\subset X\), one has
\[
  \|T^n|_M\|>\frac12\|T^n\|
\]
for all sufficiently large \(n\).  Since
\(\sum_i f(2^k\alpha_i)<\infty\), each product
\[
  \prod_i \bigl(1+f(2^k\alpha_i)\bigr)
\]
converges.  Choose integers \(m_k\uparrow\infty\) with
\(\alpha_i\le2^{-k}\) for \(i\ge m_k\) and
\[
  \prod_{i=m_k}^{\infty}\bigl(1+f(2^k\alpha_i)\bigr)\le2 .
\]
Auge's induction then constructs \(n_j\uparrow\infty\) and unit vectors
\(u_i\in X\) so that the partial sums \(x_l=\sum_{i=1}^l\alpha_i u_i\)
converge to some \(x\in X\), while
\[
  \|T^{n_j}x\|\ge \frac{\alpha_j}{2}\|T^{n_j}\|
  \qquad (j\ge1).
\]
Indeed the only estimate involving the modulus is the source inequality
\[
  \left\|\frac{s_l}{\|s_l\|}+2^k\alpha_{l+1}y\right\|
  \le (1+\beta_{l+1})\bigl(1+f(2^k\alpha_{l+1})\bigr),
\]
valid on a suitable finite-codimensional subspace; our summability of
\(f(2^k\alpha_i)\) is exactly what is needed to keep the block products
bounded and the partial sums Cauchy.

Consequently,
\[
  \sum_{n=1}^{\infty}\rho\!\left(\frac{\|T^n x\|}{\|T^n\|}\right)
  \ge
  \sum_{j=1}^{\infty}\rho(\alpha_j/2)
  =\infty,
\]
as required.

\section*{Scope}

The result is a partial positive answer to the source remark.  It removes
the dyadic regularity assumption from the AUS modulus for all dyadically
doubling comparison gauges, including the power gauges used to compute the
classical exponents.  It leaves open the genuinely irregular case where
\(\rho\) itself can decay much faster under fixed dilations.

\end{document}
